\documentclass[10pt,aspectratio=169,serif]{beamer}

\usepackage[fontset=none, scheme=plain]{ctex}
\IfFontExistsTF{Songti SC}{%
	\setCJKmainfont{Songti SC}[AutoFakeSlant=0.2]%
	\setCJKsansfont{Songti SC}[AutoFakeSlant=0.2]%
}{%
	\setCJKmainfont{FandolSong-Regular}[Extension=.otf,BoldFont=FandolSong-Bold,AutoFakeSlant=0.2]%
	\setCJKsansfont{FandolSong-Regular}[Extension=.otf,BoldFont=FandolSong-Bold]%
}
\usepackage{amsmath,amssymb,graphicx,tikz,booktabs,multicol,calligra}
\usepackage{pgfplots}
\pgfplotsset{compat=1.18}
\usepackage{tcolorbox}
\tcbuselibrary{skins}
\usepackage{XyygU}
\usepackage{caption}
\captionsetup[figure]{name=图,labelsep=period}

\usetikzlibrary{arrows.meta,positioning,calc}

\setlength{\abovedisplayskip}{0pt plus 2pt minus 2pt}
\setlength{\belowdisplayskip}{0pt plus 2pt minus 2pt}

\title[习题5.2]{习题5.2 三角函数的概念}
\subtitle{复习巩固 \& 综合运用 \& 拓广探索}
\author{}
\institute{}
\date{2026-10-16}

\begin{document}

\begin{frame}
  \titlepage
\end{frame}

% ========== 第1题 ==========
\begin{frame}[t]{第1题：求三角函数值}
\textbf{1.} 用定义法、公式一求下列角的三个三角函数值（可用计算工具）：
\begin{enumerate}
  \item $-\dfrac{17\pi}{3}$；
  \item $\dfrac{21\pi}{4}$；
  \item $-\dfrac{23\pi}{6}$；
  \item $1500^\circ$。
\end{enumerate}
\end{frame}

\begin{frame}[t]{第1题 解答}
\begin{tcolorbox}[colback=orange!10,colframe=orange,title=\textbf{解}]
\small
(1) $-\dfrac{17\pi}{3}+3\times 2\pi=-\dfrac{17\pi}{3}+\dfrac{18\pi}{3}=\dfrac{\pi}{3}$：
$\sin=\dfrac{\sqrt{3}}{2},\ \cos=\dfrac{1}{2},\ \tan=\sqrt{3}$。

(2) $\dfrac{21\pi}{4}-2\times 2\pi=\dfrac{21\pi}{4}-\dfrac{16\pi}{4}=\dfrac{5\pi}{4}$：
$\sin=-\dfrac{\sqrt{2}}{2},\ \cos=-\dfrac{\sqrt{2}}{2},\ \tan=1$。

(3) $-\dfrac{23\pi}{6}+2\times 2\pi=-\dfrac{23\pi}{6}+\dfrac{24\pi}{6}=\dfrac{\pi}{6}$：
$\sin=\dfrac{1}{2},\ \cos=\dfrac{\sqrt{3}}{2},\ \tan=\dfrac{\sqrt{3}}{3}$。

(4) $1500^\circ-4\times 360^\circ=1500^\circ-1440^\circ=60^\circ$：
$\sin=\dfrac{\sqrt{3}}{2},\ \cos=\dfrac{1}{2},\ \tan=\sqrt{3}$。
\end{tcolorbox}
\end{frame}

% ========== 第2题 ==========
\begin{frame}[t]{第2题：由终边上点求三角函数值}
\textbf{2.} 已知角 $\alpha$ 的终边上有一点 $P$ 的坐标是 $(3a,4a)$，其中 $a\neq 0$，求 $\sin\alpha$，$\cos\alpha$，$\tan\alpha$ 的值。
\end{frame}

\begin{frame}[t]{第2题 解答}
\begin{tcolorbox}[colback=orange!10,colframe=orange,title=\textbf{解}]
点 $P(3a,4a)$ 到原点距离：
\[
r = \sqrt{(3a)^2+(4a)^2} = \sqrt{25a^2} = 5|a|.
\]

\vspace{0.5em}
\textbf{当 $a>0$ 时}（$P$ 在第一象限）：
\[
\sin\alpha=\frac{4a}{5a}=\frac{4}{5},\quad
\cos\alpha=\frac{3a}{5a}=\frac{3}{5},\quad
\tan\alpha=\frac{4a}{3a}=\frac{4}{3}.
\]

\vspace{0.3em}
\textbf{当 $a<0$ 时}（$P$ 在第三象限）：
\[
\sin\alpha=-\frac{4}{5},\quad
\cos\alpha=-\frac{3}{5},\quad
\tan\alpha=\frac{4}{3}.
\]
\end{tcolorbox}
\end{frame}

% ========== 第3题 ==========
\begin{frame}[t]{第3题：计算特殊角三角函数值}
\textbf{3.} 计算：
\begin{enumerate}
  \item $6\sin(-90^\circ)+3\sin 0^\circ-8\sin 270^\circ+12\cos 180^\circ$；
  \item $10\cos 270^\circ+4\sin 0^\circ+9\tan 0^\circ+15\cos 360^\circ$；
  \item $2\cos\dfrac{\pi}{2}-\tan\dfrac{\pi}{4}+\dfrac{3}{4}\tan^2\dfrac{\pi}{6}-\sin\dfrac{\pi}{6}+\cos^2\dfrac{\pi}{6}+\sin\dfrac{3\pi}{2}$；
  \item $\sin^2\dfrac{\pi}{3}+\cos^4\dfrac{3\pi}{2}-\tan^2\dfrac{\pi}{3}$。
\end{enumerate}
\end{frame}

\begin{frame}[t]{第3题 解答}
\begin{tcolorbox}[colback=orange!10,colframe=orange,title=\textbf{解}]
\small
(1) 原式 $=6(-1)+3\cdot 0-8(-1)+12(-1)=-6+0+8-12=\alert{-10}$。

(2) 原式 $=10\cdot 0+4\cdot 0+9\cdot 0+15\cdot 1=\alert{15}$。

(3) 原式 $=2\cdot 0-1+\dfrac{3}{4}\cdot\dfrac{1}{3}-\dfrac{1}{2}+\dfrac{3}{4}+(-1)$
$=0-1+\dfrac{1}{4}-\dfrac{1}{2}+\dfrac{3}{4}-1=\alert{-1}$。

(4) 原式 $=\left(\dfrac{\sqrt{3}}{2}\right)^2+0-(\sqrt{3})^2=\dfrac{3}{4}-3=\alert{-\dfrac{9}{4}}$。
\end{tcolorbox}
\end{frame}

% ========== 第4题 ==========
\begin{frame}[t]{第4题：化简}
\textbf{4.} 化简：
\begin{enumerate}
  \item $a\sin 0^\circ+b\cos 90^\circ+c\tan 180^\circ$；
  \item $-p^2\cos 180^\circ+q^2\sin 90^\circ-2pq\cos 0^\circ$；
  \item $a^2\cos 2\pi-b^2\sin\dfrac{3\pi}{2}+ab\cos\pi-ab\sin\dfrac{\pi}{2}$；
  \item $m\tan 0+n\cos\dfrac{\pi}{2}-p\sin\pi-q\cos\dfrac{3\pi}{2}-r\sin 2\pi$。
\end{enumerate}
\end{frame}

\begin{frame}[t]{第4题 解答}
\begin{tcolorbox}[colback=orange!10,colframe=orange,title=\textbf{解}]
\small
(1) 原式 $=a\cdot 0+b\cdot 0+c\cdot 0=\alert{0}$。

(2) 原式 $=-p^2(-1)+q^2\cdot 1-2pq\cdot 1=p^2+q^2-2pq=\alert{(p-q)^2}$。

(3) 原式 $=a^2\cdot 1-b^2(-1)+ab(-1)-ab\cdot 1=a^2+b^2-ab-ab=\alert{a^2+b^2-2ab}$。

(4) 原式 $=m\cdot 0+n\cdot 0-p\cdot 0-q\cdot 0-r\cdot 0=\alert{0}$。
\end{tcolorbox}
\end{frame}

% ========== 第5题 ==========
\begin{frame}[t]{第5题：确定三角函数值的符号}
\textbf{5.} 确定下列三角函数值的符号：
\begin{enumerate}
  \item $\sin 186^\circ$；
  \item $\tan 505^\circ$；
  \item $\sin 7.6\pi$；
  \item $\tan\left(-\dfrac{23\pi}{4}\right)$；
  \item $\cos 940^\circ$；
  \item $\cos\left(-\dfrac{59\pi}{17}\right)$。
\end{enumerate}
\end{frame}

\begin{frame}[t]{第5题 解答}
\begin{tcolorbox}[colback=orange!10,colframe=orange,title=\textbf{解}]
\small
(1) $186^\circ$ 是第三象限角，$\sin 186^\circ\alert{<0}$。

(2) $505^\circ-360^\circ=145^\circ$，第二象限，$\tan 505^\circ\alert{<0}$。

(3) $7.6\pi-3\times 2\pi=7.6\pi-6\pi=1.6\pi=\dfrac{8\pi}{5}$，第四象限，$\sin 7.6\pi\alert{<0}$。

(4) $-\dfrac{23\pi}{4}+3\times 2\pi=-\dfrac{23\pi}{4}+\dfrac{24\pi}{4}=\dfrac{\pi}{4}$，第一象限，$\tan\alert{>0}$。

(5) $940^\circ-2\times 360^\circ=220^\circ$，第三象限，$\cos 940^\circ\alert{<0}$。

(6) $-\dfrac{59\pi}{17}+4\times 2\pi=-\dfrac{59\pi}{17}+\dfrac{136\pi}{17}=\dfrac{77\pi}{17}=4\pi+\dfrac{9\pi}{17}$，$\dfrac{9\pi}{17}$ 是第二象限角，$\cos\alert{<0}$。
\end{tcolorbox}
\end{frame}

% ========== 第6题 ==========
\begin{frame}[t]{第6题：已知一个三角函数值求其他}
\textbf{6.}
\begin{enumerate}
  \item 已知 $\sin\alpha=-\dfrac{\sqrt{3}}{2}$，且 $\alpha$ 为第四象限角，求 $\cos\alpha$，$\tan\alpha$ 的值；
  \item 已知 $\cos\alpha=-\dfrac{5}{13}$，且 $\alpha$ 为第二象限角，求 $\sin\alpha$，$\tan\alpha$ 的值；
  \item 已知 $\tan\alpha=-\dfrac{3}{4}$，求 $\sin\alpha$，$\cos\alpha$ 的值；
  \item 已知 $\cos\alpha=0.68$，求 $\sin\alpha$，$\tan\alpha$ 的值（精确到 $0.01$）。
\end{enumerate}
\end{frame}

\begin{frame}[t]{第6题 解答}
\begin{tcolorbox}[colback=orange!10,colframe=orange,title=\textbf{解}]
\small
(1) 第四象限：$\cos\alpha=\sqrt{1-\dfrac{3}{4}}=\dfrac{1}{2}$，$\tan\alpha=\dfrac{-\sqrt{3}/2}{1/2}=-\sqrt{3}$。

(2) 第二象限：$\sin\alpha=\sqrt{1-\dfrac{25}{169}}=\dfrac{12}{13}$，$\tan\alpha=\dfrac{12/13}{-5/13}=-\dfrac{12}{5}$。

(3) $\tan\alpha=-\dfrac{3}{4}<0$，$\alpha$ 在第二或第四象限。
由 $1+\tan^2\alpha=\dfrac{1}{\cos^2\alpha}$，得 $\cos^2\alpha=\dfrac{16}{25}$。
第二象限：$\cos\alpha=-\dfrac{4}{5},\ \sin\alpha=\dfrac{3}{5}$；
第四象限：$\cos\alpha=\dfrac{4}{5},\ \sin\alpha=-\dfrac{3}{5}$。

(4) $\sin\alpha=\pm\sqrt{1-0.68^2}=\pm\sqrt{0.5376}\approx\pm 0.73$。
$\cos\alpha=0.68>0$，$\alpha$ 在第一或第四象限：
第一象限 $\sin\approx 0.73,\ \tan\approx 1.07$；
第四象限 $\sin\approx -0.73,\ \tan\approx -1.07$。
\end{tcolorbox}
\end{frame}

% ========== 第7题 ==========
\begin{frame}[t]{第7题：求函数值}
\textbf{7.} 根据下列条件求函数 $f(x)=\sin\left(x+\dfrac{\pi}{4}\right)+2\sin\left(x-\dfrac{\pi}{4}\right)-4\cos 2x+3\cos\left(x+\dfrac{3\pi}{4}\right)$ 的值：
\begin{enumerate}
  \item $x=\dfrac{\pi}{4}$；
  \item $x=\dfrac{3\pi}{4}$。
\end{enumerate}
\end{frame}

\begin{frame}[t]{第7题 解答}
\begin{tcolorbox}[colback=orange!10,colframe=orange,title=\textbf{解}]
(1) $x=\dfrac{\pi}{4}$：
\[
f\left(\frac{\pi}{4}\right)=\sin\frac{\pi}{2}+2\sin 0-4\cos\frac{\pi}{2}+3\cos\pi
=1+0-0+3(-1)=\alert{-2}.
\]

(2) $x=\dfrac{3\pi}{4}$：
\[
f\left(\frac{3\pi}{4}\right)=\sin\pi+2\sin\frac{\pi}{2}-4\cos\frac{3\pi}{2}+3\cos\frac{3\pi}{2}
=0+2\cdot 1-4\cdot 0+3\cdot 0=\alert{2}.
\]
\end{tcolorbox}
\end{frame}

% ========== 第8题 ==========
\begin{frame}[t]{第8题：确定式子的符号}
\textbf{8.} 确定下列式子的符号：
\begin{enumerate}
  \item $\tan 125^\circ\sin 273^\circ$；
  \item $\dfrac{\tan 108^\circ}{\cos 305^\circ}$；
  \item $\sin\dfrac{5\pi}{4}\cos\dfrac{4\pi}{5}\tan\dfrac{11\pi}{6}$；
  \item $\dfrac{\cos\dfrac{5\pi}{6}\tan\dfrac{11\pi}{6}}{\sin\dfrac{2\pi}{3}}$。
\end{enumerate}
\end{frame}

\begin{frame}[t]{第8题 解答}
\begin{tcolorbox}[colback=orange!10,colframe=orange,title=\textbf{解}]
\small
(1) $125^\circ$ 第二象限 $\tan<0$；$273^\circ$ 第四象限 $\sin<0$；乘积 $\alert{>0}$。

(2) $108^\circ$ 第二象限 $\tan<0$；$305^\circ$ 第四象限 $\cos>0$；商 $\alert{<0}$。

(3) $\dfrac{5\pi}{4}$ 第三象限 $\sin<0$；$\dfrac{4\pi}{5}$ 第二象限 $\cos<0$；$\dfrac{11\pi}{6}$ 第四象限 $\tan<0$；三者乘积 $(-)(-)(-)=\alert{<0}$。

(4) $\dfrac{5\pi}{6}$ 第二象限 $\cos<0$；$\dfrac{11\pi}{6}$ 第四象限 $\tan<0$；$\dfrac{2\pi}{3}$ 第二象限 $\sin>0$；$\dfrac{(-)(-)}{(+)}=\alert{>0}$。
\end{tcolorbox}
\end{frame}

% ========== 第9题 ==========
\begin{frame}[t]{第9题：求三角函数值}
\textbf{9.} 求下列三角函数值（可用计算工具，第(1)(3)(4)题精确到 $0.0001$）：
\begin{enumerate}
  \item $\sin\left(-\dfrac{67\pi}{12}\right)$；
  \item $\tan\left(-\dfrac{15\pi}{4}\right)$；
  \item $\cos 398^\circ 13'$；
  \item $\tan 766^\circ 15'$。
\end{enumerate}
\end{frame}

\begin{frame}[t]{第9题 解答}
\begin{tcolorbox}[colback=orange!10,colframe=orange,title=\textbf{解}]
\small
(1) $-\dfrac{67\pi}{12}+6\times 2\pi=-\dfrac{67\pi}{12}+\dfrac{144\pi}{12}=\dfrac{77\pi}{12}=6\pi+\dfrac{5\pi}{12}$，
$\sin\left(-\dfrac{67\pi}{12}\right)=\sin\dfrac{5\pi}{12}=\sin 75^\circ\approx 0.9659$。

(2) $-\dfrac{15\pi}{4}+2\times 2\pi=-\dfrac{15\pi}{4}+\dfrac{16\pi}{4}=\dfrac{\pi}{4}$，
$\tan\left(-\dfrac{15\pi}{4}\right)=\tan\dfrac{\pi}{4}=1$。

(3) $398^\circ 13'-360^\circ=38^\circ 13'$，
$\cos 398^\circ 13'=\cos 38^\circ 13'\approx 0.7854$。

(4) $766^\circ 15'-2\times 360^\circ=46^\circ 15'$，
$\tan 766^\circ 15'=\tan 46^\circ 15'\approx 1.0445$。
\end{tcolorbox}
\end{frame}

% ========== 第10题 ==========
\begin{frame}[t]{第10题：充要条件证明}
\textbf{10.} 求证：
\begin{enumerate}
  \item 角 $\theta$ 为第二或第三象限角的充要条件是 $\sin\theta\tan\theta<0$；
  \item 角 $\theta$ 为第三或第四象限角的充要条件是 $\cos\theta\tan\theta<0$；
  \item 角 $\theta$ 为第一或第四象限角的充要条件是 $\dfrac{\sin\theta}{\tan\theta}>0$；
  \item 角 $\theta$ 为第一或第三象限角的充要条件是 $\sin\theta\cos\theta>0$。
\end{enumerate}
\end{frame}

\begin{frame}[t]{第10题 解答}
\begin{tcolorbox}[colback=orange!10,colframe=orange,title=\textbf{证明}]
\small
由各象限符号：I($\sin+, \cos+, \tan+$)，II($\sin+, \cos-, \tan-$)，III($\sin-, \cos-, \tan+$)，IV($\sin-, \cos+, \tan-$)。

(1) $\sin\theta\tan\theta<0$ 要求 $\sin$ 与 $\tan$ 异号：II($+\cdot-<0$) 或 III($-\cdot+<0$)。$\checkmark$

(2) $\cos\theta\tan\theta<0$ 要求 $\cos$ 与 $\tan$ 异号：III($-\cdot+<0$) 或 IV($+\cdot-<0$)。$\checkmark$

(3) $\dfrac{\sin\theta}{\tan\theta}=\cos\theta>0$，要求 $\cos>0$：I 或 IV。$\checkmark$

(4) $\sin\theta\cos\theta>0$ 要求同号：I($+\cdot+$) 或 III($-\cdot-$)。$\checkmark$
\end{tcolorbox}
\end{frame}

% ========== 第11题 ==========
\begin{frame}[t]{第11题：已知正弦求余弦正切}
\textbf{11.} 已知 $\sin x=-\dfrac{1}{3}$，求 $\cos x$，$\tan x$ 的值。
\end{frame}

\begin{frame}[t]{第11题 解答}
\begin{tcolorbox}[colback=orange!10,colframe=orange,title=\textbf{解}]
\[
\cos^2 x = 1-\sin^2 x = 1-\frac{1}{9} = \frac{8}{9},
\]
\[
\cos x = \pm\frac{2\sqrt{2}}{3}.
\]

\vspace{0.5em}
$\sin x=-\dfrac{1}{3}<0$，$x$ 在第三或第四象限：

\vspace{0.3em}
\textbf{第三象限：} $\cos x=-\dfrac{2\sqrt{2}}{3}$，
$\tan x=\dfrac{-1/3}{-2\sqrt{2}/3}=\dfrac{\sqrt{2}}{4}$；

\vspace{0.3em}
\textbf{第四象限：} $\cos x=\dfrac{2\sqrt{2}}{3}$，
$\tan x=\dfrac{-1/3}{2\sqrt{2}/3}=-\dfrac{\sqrt{2}}{4}$。
\end{tcolorbox}
\end{frame}

% ========== 第12题 ==========
\begin{frame}[t]{第12题：已知正切求三角式}
\textbf{12.} 已知 $\tan\alpha=\sqrt{3}$，$\pi<\alpha<\dfrac{3\pi}{2}$，求 $\cos\alpha-\sin\alpha$ 的值。
\end{frame}

\begin{frame}[t]{第12题 解答}
\begin{tcolorbox}[colback=orange!10,colframe=orange,title=\textbf{解}]
$\tan\alpha=\sqrt{3}$ 且 $\pi<\alpha<\dfrac{3\pi}{2}$，所以 $\alpha=\dfrac{4\pi}{3}$。

\vspace{0.5em}
\[
\cos\alpha=\cos\frac{4\pi}{3}=-\frac{1}{2},\quad
\sin\alpha=\sin\frac{4\pi}{3}=-\frac{\sqrt{3}}{2}.
\]
\[
\cos\alpha-\sin\alpha = -\frac{1}{2}-\left(-\frac{\sqrt{3}}{2}\right)
= \frac{\sqrt{3}-1}{2}.
\]
\end{tcolorbox}
\end{frame}

% ========== 第13题 ==========
\begin{frame}[t]{第13题：用一个三角函数表示另两个}
\textbf{13.} 已知角 $\alpha$ 的终边不在坐标轴上，
\begin{enumerate}
  \item 用 $\cos\alpha$ 表示 $\sin\alpha$，$\tan\alpha$；
  \item 用 $\sin\alpha$ 表示 $\cos\alpha$，$\tan\alpha$。
\end{enumerate}
\end{frame}

\begin{frame}[t]{第13题 解答}
\begin{tcolorbox}[colback=orange!10,colframe=orange,title=\textbf{解}]
(1) 由 $\sin^2\alpha+\cos^2\alpha=1$：
\[
\sin\alpha=\pm\sqrt{1-\cos^2\alpha},\quad
\tan\alpha=\frac{\sin\alpha}{\cos\alpha}=\pm\frac{\sqrt{1-\cos^2\alpha}}{\cos\alpha}.
\]
（符号由 $\alpha$ 所在象限决定）

(2) 同理：
\[
\cos\alpha=\pm\sqrt{1-\sin^2\alpha},\quad
\tan\alpha=\frac{\sin\alpha}{\cos\alpha}=\pm\frac{\sin\alpha}{\sqrt{1-\sin^2\alpha}}.
\]
\end{tcolorbox}
\end{frame}

% ========== 第14题 ==========
\begin{frame}[t]{第14题：证明恒等式}
\textbf{14.} 求证：
\begin{enumerate}
  \item $\dfrac{1-2\sin x\cos x}{\cos^2 x-\sin^2 x}=\dfrac{1-\tan x}{1+\tan x}$；
  \item $\tan^2\alpha-\sin^2\alpha=\tan^2\alpha\sin^2\alpha$；
  \item $(\cos\beta-1)^2+\sin^2\beta=2-2\cos\beta$；
  \item $\sin^4 x+\cos^4 x=1-2\sin^2 x\cos^2 x$。
\end{enumerate}
\end{frame}

\begin{frame}[t]{第14题 解答}
\begin{tcolorbox}[colback=orange!10,colframe=orange,title=\textbf{证明}]
\small
(1) 左边分子 $1-2\sin x\cos x=(\cos x-\sin x)^2$，分母 $\cos^2 x-\sin^2 x=(\cos x-\sin x)(\cos x+\sin x)$，约去 $\cos x-\sin x$，得 $\dfrac{\cos x-\sin x}{\cos x+\sin x}=\dfrac{1-\tan x}{1+\tan x}$。$\checkmark$

(2) 右边 $\tan^2\alpha\sin^2\alpha=\tan^2\alpha(1-\cos^2\alpha)=\tan^2\alpha-\sin^2\alpha=$ 左边。$\checkmark$

(3) 左边 $=\cos^2\alpha-2\cos\alpha+1+\sin^2\alpha=(\cos^2\alpha+\sin^2\alpha)+1-2\cos\alpha=2-2\cos\alpha$。$\checkmark$

(4) 右边 $=1-2\sin^2 x\cos^2 x=(\sin^2 x+\cos^2 x)^2-2\sin^2 x\cos^2 x=\sin^4 x+\cos^4 x=$ 左边。$\checkmark$
\end{tcolorbox}
\end{frame}

% ========== 第15题 ==========
\begin{frame}[t]{第15题：齐次式求值}
\textbf{15.} 已知 $\tan\alpha=2$，求 $\dfrac{\sin\alpha+\cos\alpha}{\sin\alpha-\cos\alpha}$ 的值。
\end{frame}

\begin{frame}[t]{第15题 解答}
\begin{tcolorbox}[colback=orange!10,colframe=orange,title=\textbf{解}]
分子分母同除以 $\cos\alpha$：
\[
\frac{\sin\alpha+\cos\alpha}{\sin\alpha-\cos\alpha}
= \frac{\dfrac{\sin\alpha}{\cos\alpha}+1}{\dfrac{\sin\alpha}{\cos\alpha}-1}
= \frac{\tan\alpha+1}{\tan\alpha-1}.
\]
代入 $\tan\alpha=2$：
\[
= \frac{2+1}{2-1} = \alert{3}.
\]
\end{tcolorbox}
\end{frame}

% ========== 第16题 ==========
\begin{frame}[t]{第16题：化简带根号的式子}
\textbf{16.} 化简 $\sqrt{\dfrac{1+\sin\alpha}{1-\sin\alpha}}-\sqrt{\dfrac{1-\sin\alpha}{1+\sin\alpha}}$，其中 $\alpha$ 为第二象限角。
\end{frame}

\begin{frame}[t]{第16题 解答}
\begin{tcolorbox}[colback=orange!10,colframe=orange,title=\textbf{解}]
$\alpha$ 为第二象限角，$\cos\alpha<0$。
对两式分别有理化：
\[
\sqrt{\frac{1+\sin\alpha}{1-\sin\alpha}}
= \sqrt{\frac{(1+\sin\alpha)^2}{1-\sin^2\alpha}}
= \sqrt{\frac{(1+\sin\alpha)^2}{\cos^2\alpha}}
= \frac{1+\sin\alpha}{|\cos\alpha|}.
\]
因为 $\cos\alpha<0$，$|\cos\alpha|=-\cos\alpha$：
\[
= \frac{1+\sin\alpha}{-\cos\alpha}.
\]
同理 $\sqrt{\dfrac{1-\sin\alpha}{1+\sin\alpha}}=\dfrac{1-\sin\alpha}{-\cos\alpha}$。

原式 $=\dfrac{1+\sin\alpha}{\!-\cos\alpha}-\dfrac{1-\sin\alpha}{\!-\cos\alpha}=\dfrac{2\sin\alpha}{-\cos\alpha}=\alert{-2\tan\alpha}$。
\end{tcolorbox}
\end{frame}

% ========== 第17题 ==========
\begin{frame}[t]{第17题：推导更多关系式}
\textbf{17.} 从本节的例 7 可以看出，$\dfrac{\cos x}{1-\sin x}=\dfrac{1+\sin x}{\cos x}$ 就是 $\sin^2 x+\cos^2 x=1$ 的一个变形。你能利用同角三角函数的基本关系推导出更多的关系式吗？
\end{frame}

\begin{frame}[t]{第17题 解答}
\begin{tcolorbox}[colback=orange!10,colframe=orange,title=\textbf{参考答案}]
由 $\sin^2 x+\cos^2 x=1$，可推出：
\begin{itemize}
  \item $\sin^2 x = 1-\cos^2 x=(1-\cos x)(1+\cos x)$；
  \item $\cos^2 x = 1-\sin^2 x=(1-\sin x)(1+\sin x)$；
  \item $1 = \sin^2 x+\cos^2 x$；
  \item $(\sin x+\cos x)^2 = 1+2\sin x\cos x$；
  \item $(\sin x-\cos x)^2 = 1-2\sin x\cos x$；
  \item $\dfrac{1-\cos x}{\sin x}=\dfrac{\sin x}{1+\cos x}$；
  \item $1+\tan^2 x=\dfrac{1}{\cos^2 x}$，$1+\cot^2 x=\dfrac{1}{\sin^2 x}$。
\end{itemize}
\end{tcolorbox}
\end{frame}

% ========== 第18题 ==========
\begin{frame}[t]{第18题：探究恒等式}
\textbf{18.}
\begin{enumerate}
  \item 分别计算 $\sin^4\dfrac{\pi}{3}-\cos^4\dfrac{\pi}{3}$ 和 $\sin^2\dfrac{\pi}{3}-\cos^2\dfrac{\pi}{3}$ 的值，你有什么发现？
  \item 任取一个 $\alpha$ 的值，分别计算 $\sin^4\alpha-\cos^4\alpha$，$\sin^2\alpha-\cos^2\alpha$，你又有什么发现？
  \item 证明：$\forall x\in\mathbf{R}$，$\sin^2 x-\cos^2 x=\sin^4 x-\cos^4 x$。
\end{enumerate}
\end{frame}

\begin{frame}[t]{第18题 解答}
\begin{tcolorbox}[colback=orange!10,colframe=orange,title=\textbf{解}]
(1) $\sin\dfrac{\pi}{3}=\dfrac{\sqrt{3}}{2}$，$\cos\dfrac{\pi}{3}=\dfrac{1}{2}$：
\[
\sin^4\frac{\pi}{3}-\cos^4\frac{\pi}{3}=\frac{9}{16}-\frac{1}{16}=\frac{1}{2},
\]
\[
\sin^2\frac{\pi}{3}-\cos^2\frac{\pi}{3}=\frac{3}{4}-\frac{1}{4}=\frac{1}{2}.
\]
两者相等。

(2) 任取 $\alpha$，两者均相等。

(3) 证明：
\[
\sin^4 x-\cos^4 x=(\sin^2 x-\cos^2 x)(\sin^2 x+\cos^2 x).
\]
由 $\sin^2 x+\cos^2 x=1$，得 $\sin^4 x-\cos^4 x=\sin^2 x-\cos^2 x$。$\checkmark$
\end{tcolorbox}
\end{frame}

\begin{frame}[c]
  \begin{center}
    {\Huge\calligra Thanks!}
  \end{center}
\end{frame}

\end{document}
